% Figure: Brayton cycle with regeneration on temperature-entropy axes. The
% turbine exhaust (state 4) preheats the compressed air (state 2) toward state x
% before it reaches the combustor, cutting the fuel heat.
\begin{figure}[htbp]
\centering
\begin{tikzpicture}
\begin{axis}[
  width=12cm, height=8cm,
  xlabel={specific entropy $s$ (arb.)},
  ylabel={temperature $T$ (K)},
  xmin=0, xmax=6, ymin=250, ymax=1550,
  axis lines=left,
  tick label style={font=\scriptsize},
  label style={font=\small},
  clip=false,
]
% compression 1->2 (nearly vertical, entropy rise small for near-isentropic)
\addplot[engblue, very thick] coordinates {(1.0,290) (1.15,643)};
% preheat in regenerator 2->x (constant pressure, entropy rises)
\addplot[engamber, very thick, smooth] coordinates {(1.15,643) (2.4,760)};
% combustor x->3 (constant pressure)
\addplot[engred, very thick, smooth] coordinates {(2.4,760) (4.6,1450)};
% turbine 3->4 (near vertical down)
\addplot[engorange, very thick] coordinates {(4.6,1450) (4.75,801)};
% regenerator cooling of exhaust 4->y (constant pressure down)
\addplot[enggray, very thick, smooth] coordinates {(4.75,801) (3.5,690)};
% exhaust reject y->1 (down to inlet, constant pressure)
\addplot[enggray, very thick, dashed, smooth] coordinates {(3.5,690) (1.0,290)};
% markers
\filldraw[engblack] (axis cs:1.0,290) circle [radius=1.6pt];
\node[engblack, font=\scriptsize, anchor=north east] at (axis cs:1.0,290) {1};
\filldraw[engblack] (axis cs:1.15,643) circle [radius=1.6pt];
\node[engblack, font=\scriptsize, anchor=east] at (axis cs:1.15,643) {2};
\filldraw[engblack] (axis cs:2.4,760) circle [radius=1.6pt];
\node[engblack, font=\scriptsize, anchor=north] at (axis cs:2.4,760) {x};
\filldraw[engblack] (axis cs:4.6,1450) circle [radius=1.6pt];
\node[engblack, font=\scriptsize, anchor=south east] at (axis cs:4.6,1450) {3};
\filldraw[engblack] (axis cs:4.75,801) circle [radius=1.6pt];
\node[engblack, font=\scriptsize, anchor=west] at (axis cs:4.75,801) {4};
\end{axis}
\end{tikzpicture}
\caption[Regenerated Brayton cycle]{Adding a regenerator to the Brayton cycle
recovers part of the hot turbine exhaust. The exhaust at state 4 flows through a
counterflow exchanger and preheats the compressed air from state 2 to state x
before it enters the combustor, so the fuel need only carry the gas from x to 3
rather than from 2 to 3. The net work is unchanged, but the heat added falls, so
the efficiency rises. Regeneration pays only when the turbine exhaust is hotter
than the compressor exit, which for a fixed turbine inlet temperature means at low
pressure ratios; at high pressure ratios the compressor exit overtakes the exhaust
and the regenerator would run backward, so it is omitted.}
\label{fig:vol04ch01:regeneration-ts}
\end{figure}
